% Generated draft: requires mathematical and source review.
\documentclass[11pt]{article}
\usepackage[margin=25mm]{geometry}
\usepackage{fontspec}
\setmainfont{Latin Modern Roman}
\usepackage{amsmath,amssymb,mathtools,unicode-math}
\setmathfont{Latin Modern Math}
\usepackage{enumitem,needspace,xurl,hyperref}
\setlength{\parindent}{0pt}
\setlength{\parskip}{5pt}
\setlength{\emergencystretch}{3em}
\newcommand{\N}{\mathbb{N}}
\newcommand{\Z}{\mathbb{Z}}
\newcommand{\Q}{\mathbb{Q}}
\newcommand{\R}{\mathbb{R}}
\newcommand{\C}{\mathbb{C}}
\newcommand{\F}{\mathbb{F}}
\newcommand{\A}{\mathbb{A}}
\newcommand{\PP}{\mathbb{P}}
\newcommand{\E}{\mathbb{E}}
\newcommand{\eps}{\varepsilon}
\newcommand{\dd}{\,\mathrm{d}}
\begin{document}
\section*{UnsolvedMath ID: 30002672}
\textbf{Source problem code:} catalogue entry
\subsection{Definitions and mathematical statement}
\textbf{Two independent questions.} The supplied record combines a line-discrepancy question and an orthic-triangle covering question. They are stated separately below.

\textbf{(1) Line-discrepancy centers.} Let \(\mathcal L\) be the set of affine lines in \(\R^2\). A coloring of a finite set \(P\subset\R^2\) with two available colors is represented by a map \(c:P\to\{-1,1\}\). Write \(D_{P,c}(\ell)\) for the line discrepancy mentioned in the source. The supplied record does not define this quantity; in particular, it does not specify which points or regions associated with \(\ell\) enter its computation. Thus the following is a conditional formulation whose meaning depends on recovering that definition.

Is it true that for every real \(\alpha>0\) there exists a real \(\beta<\alpha\) such that, for every positive integer \(n\), every set \(P\subset\R^2\) with \(|P|=n\), and every coloring \(c:P\to\{-1,1\}\),
\[
\bigl[D_{P,c}(\ell)\leq\alpha n\ \text{for every }\ell\in\mathcal L\bigr]
\quad\Longrightarrow\quad
\bigl[\exists q\in\R^2\ \forall\ell\in\mathcal L,
\ q\in\ell\Longrightarrow D_{P,c}(\ell)\leq\beta n\bigr]?
\]
Here \(\beta\) is chosen before \(n\), \(P\), and \(c\), while \(q\) may depend on the colored point set. This reads the source as a uniform assertion over positive integers \(n\), and reads \(q\) as an arbitrary point of the ambient plane; these conventions require confirmation against the report. The source states only \(\beta<\alpha\), with no explicit lower bound on \(\beta\). The coloring notation allows either color class to be empty; whether both colors must occur is unspecified.

\textbf{(2) Orthic-triangle coverings.} For a nondegenerate acute Euclidean triangle \(U=ABC\), let \(H_A,H_B,H_C\) be the perpendicular projections of \(A,B,C\) onto the lines \(BC,CA,AB\), respectively. Its orthic triangle is \(H_AH_BH_C\), and define
\[
\xi(U)=\lVert H_A-H_B\rVert+\lVert H_B-H_C\rVert+\lVert H_C-H_A\rVert.
\]
Interpreting triangles as their closed filled regions in \(\R^2\), is it true that, for every positive integer \(N\) and every collection of nondegenerate acute triangles \(T,S_1,\ldots,S_N\),
\[
T\subseteq\bigcup_{i=1}^{N}S_i
\quad\Longrightarrow\quad
\xi(T)\leq\sum_{i=1}^{N}\xi(S_i)?
\]
This interpretation permits overlaps and does not require \(S_i\subseteq T\). The supplied record does not expressly specify the boundary convention for a covering or impose containment or disjointness conditions. The closed-region interpretation therefore remains a convention requiring source confirmation.
\subsection{Short English statement}
The first question asks whether a uniform discrepancy bound for all lines can always be strictly improved for the lines through some point, with the improvement depending only on \(\alpha\). Its discrepancy function is undefined in the supplied record. The second asks whether covering an acute triangle by finitely many acute triangles forces the sum of their orthic perimeters to be at least its own. Both formulations require human mathematical review.
\subsection{Sources}
\begin{itemize}
\item Local UnsolvedMath registry, permanent ID \(30002672\), source problem code \texttt{catalogue entry}, title: \emph{Line Discrepancy Centers and Orthic-Triangle Coverings}. The supplied collection is \emph{Oberwolfach Reports Open Problems}. No individual proposer is identified; the supplied proposed year is 2014.
\item Supplied report citation: \emph{Discrete Geometry (2015)}, Oberwolfach Reports, pp. 2235--2296, DOI \texttt{10.4171/owr/2014/40}, \url{https://doi.org/10.4171/owr/2014/40}. The record describes Oberwolfach Reports as documenting mathematical workshops and their open-problem discussions. The citation's title year 2015 and the record's proposed year 2014 are retained without reconciliation.
\item Supplied literature source, labeled ``Oberwolfach report'': \url{https://ems.press/content/serial-article-files/46530}.
\item Literature assessment dated exactly \texttt{2026-08-21}: the supplied curation labels the record open and states that it combines two distinct questions, a line-discrepancy assertion and an orthic-triangle covering inequality, with no authoritative source located resolving either exact statement. This is supplied dated curation, not an independently verified literature assessment; time-sensitive status and cited claims require independent checking.
\item The supplied statement status is \texttt{corrected\_verified}. Its verification note reports that the official Oberwolfach passage was checked directly and that the question, definitions, and notation were restored without adding a stronger claim. That is a claim of the supplied registry, not independent verification by this draft.
\end{itemize}
\end{document}
